\scanpage{13}
\mspage{- 102 -}
% archive sheet number: 12
% lines 1–4
D'autre part, si $Z_{\unsure{o|0}}\handwritten{\subset} Z_1\handwritten{\subset}\ldots\handwritten{\subset}Z_n$ est une chaîne de parties fermées irréductibles de $X$ et si $x\handwritten{\in}Z_o$, pour tout \unsure{o|O}uvert $U\handwritten{\subset}X$ contenant $x$, $U\handwritten{\cap}Z_i$ est irréductible et fermé dans $U$, et comme $\overline{U\handwritten{\cap}Z}_i=Z_i$, les $U\handwritten{\cap}Z_i$ sont distincts ; donc $\dim(U)\handwritten{\geq}n$, ce qui démontre (ii).

% lines 5–6
\uline{Corollaire 1 .- Si} $(U_{\handwritten{\alpha}})$ \uline{est un recouvrement ouvert de} $X$, \uline{on a} $\dim(\handwritten{X)}=\sup_{\handwritten{\alpha}}(\dim(U_{\handwritten{\alpha}}))$.

% lines 7–8
\uline{Corollaire 2 .- Soit} $X$ \uline{un espace noethérien vérifiant l'axiome} $(T_o)$ \uline{et soit} $F$ \uline{l'ensemble des points fermés de $X$.Alors} $\dim(X)=\sup_{x\handwritten{\in}F}(\dim_x(X)\handwritten{)}$

% lines 9–10
Avec les notations de la démonstration de la prop.2 (ii), il suffi\handwritten{t} d'observer qu'il existe dans $Z_o$ un point fermé (\S{} 4,n°1,lemme 1).

% lines 11
\uline{Corollaire 3 .- Soit} $Y$ \uline{une\deleted{\unsure{s}} partie fermée de} $X$. \uline{Alors}

% lines 12
\[
 \dim(Y)+\operatorname{codim}_X(Y)\handwritten{\leq}\dim(X).\tag*{(1)}
\]

% lines 13–16
Toute chaîne de parties fermées irréductibles de $Y$ étant contenue dans une composante irréductible de $Y$, on est ramené (par la déf.2) au cas où $Y$ est irréductible, et (1) est alors conséquence des définitions 1 et 2 .

% lines 17–18
\uline{Corollaire 4 .- Soient} $Y,Z,T$ \uline{trois parties fermées irréductibles de} $X$ \uline{telles que} $Y\handwritten{\subset}Z\handwritten{\subset}T$. \uline{Alors}

% lines 19
\[
 \operatorname{codim}_T(Y)\handwritten{\geq}\operatorname{codim}_Z(Y)+\operatorname{codim}_T(Z).\tag*{(2)}
\]

% lines 20
C'est une conséquence immédiate de la déf.2 .

% lines 21–24
\uline{Proposition 3 .- Soit} $X$ \uline{un espace noethérien vérifiant l'axiome} $(T_o)$ \uline{Pour que} $\dim(X)=0$, \uline{il faut et il suffit que} $X$ \uline{soit fini et discret}. \uline{Pour qu'une partie fermée} $Y$ \uline{de} $X$ \uline{soit de codimension nulle dans} $X$, \uline{il faut et il suffit que} $Y$ \uline{contienne une composante irréductible de}

% lines 25–29
La seconde assertion est immédiate (et indépendante de toute hypothèse sur $X$). De même, si $X$ est séparé (et a fortiori discret), il est clair que $\dim(X)=0$, les seuls ensembles fermés irréductibles étant les points. Inversement, si $X$ est noethérien et satisfait à $(T_o)$, toute composante irréductible de $X$ contie\unsure{n}t un point fermé \inserted{\handwritten{(\S{} 4, n°1, lemme 1)}},
